\subsection{誤差の評価}

$a$と$x$の間で$|f^{(n+1)}(t)|\leq M$なら,積分の絶対値の評価から
\[
 |R_n(x)|\leq\frac{M}{n!}\int_{\min(a,x)}^{\max(a,x)}|x-t|^n\dd t
 =\frac{M|x-a|^{n+1}}{(n+1)!}
\]
を得る.多項式をどこまで計算すれば十分かは,この上限と,求められた精度を比較して決められる.\,$x$が$a$に近いほど高い次数の項が小さくなるだけでなく,最後に残る項も具体的な不等式で押さえられる.

なお,\,$x>a$では$(x-t)^n\geq0$である.したがって$f^{(n+1)}$の符号が分かれば,剰余項の符号も分かる.誤差の絶対値だけでなく,近似値が真の値より大きいか小さいかを判断できることがある.

\begin{example}{対数の近似}
$0\leq x\leq\frac12$の範囲で$\log(1+x)$を三次式で近似し,誤差を評価せよ.
\end{example}

中心を$a=0$に選ぶ.\,$\log(1+x)$は$x=0$で値が$0$になり,各階の微分の値も簡単だからである.微分すると
\[
 f'(x)=\frac{1}{1+x},\quad
 f''(x)=-\frac{1}{(1+x)^2},\quad
 f'''(x)=\frac{2}{(1+x)^3},\quad
 f^{(4)}(x)=-\frac{6}{(1+x)^4}
\]
なので
\[
 \log(1+x)=x-\frac{x^2}{2}+\frac{x^3}{3}+R_3(x).
\]
$0\leq t\leq x$では$-6\leq f^{(4)}(t)<0$である.したがって
\[
 -\frac{x^4}{4}\leq R_3(x)\leq0,
\]
すなわち
\[
 x-\frac{x^2}{2}+\frac{x^3}{3}-\frac{x^4}{4}
 \leq\log(1+x)
 \leq x-\frac{x^2}{2}+\frac{x^3}{3}
\]
となる.例えば$x=\frac15$なら
\[
 \frac{137}{750}-\frac{1}{2500}
 \leq\log\frac65\leq\frac{137}{750}.
\]
三次式の値を計算してから,残りの最大値を差し引けば,真の値を含む区間が得られる.

\begin{example}{定積分の近似}
$A=\int_0^1\e^{-x^2}\dd x$を有理数で上下から評価せよ.
\end{example}

$\e^{-x^2}$を直接何回も微分すると,多項式と指数関数の積が増えていく.ここでは,まず一変数の関数$\e^{-u}$を$u=0$で展開し,最後に$u=x^2$を代入する.\,$0\leq u\leq1$に対して
\[
 \e^{-u}=1-u+\frac{u^2}{2}-\frac{u^3}{6}+r_3(u).
\]
四階微分は$\e^{-u}$であり,この区間で$0<\e^{-u}\leq1$だから
\[
 0\leq r_3(u)\leq\frac{u^4}{24}.
\]
$u=x^2$を代入して積分すると
\[
\begin{aligned}
 A&=\int_0^1\left(1-x^2+\frac{x^4}{2}-\frac{x^6}{6}\right)\dd x
       +\int_0^1r_3(x^2)\dd x\\
  &=1-\frac13+\frac1{10}-\frac1{42}
       +\int_0^1r_3(x^2)\dd x.
\end{aligned}
\]
ここで
\[
 0\leq\int_0^1r_3(x^2)\dd x
 \leq\frac1{24}\int_0^1x^8\dd x=\frac1{216}
\]
なので
\[
 \boxed{\frac{26}{35}\leq A\leq\frac{26}{35}+\frac1{216}}.
\]
積分区間全体で成立する誤差評価を用意したので,その不等式を積分できた.何点かで近似がよいことを確かめるだけでは,この結論にはならない.


